\documentclass{article}
\usepackage{array}
\usepackage{enumitem}
\usepackage{float}
\usepackage{graphicx}
\usepackage{ctex}
\usepackage{indentfirst}

\newcolumntype{L}[1]{>{\raggedright\arraybackslash}p{#1}}



\title{fwwb\_qiye}
\author{子翔 周}
\date{April 2026}
\begin{document}

\maketitle

\newpage

\tableofcontents

\newpage

\section{引言}

\subsection{编写目的}

本文档作为基于OpenHarmony的智慧农业控制系统的服务外包大赛企业提交材料，严格对照企业文档相关要求编写，核心目的有五点：

一是全面呈现项目需求分析成果，清晰梳理用户需求、功能需求及稳定性、响应速度等非功能需求，确保需求逻辑清晰、覆盖全面，为项目开发奠定基础；

二是系统阐述系统设计细节，包含硬件系统架构设计、软件系统架构设计（含模块划分、接口定义等概要设计与详细设计），明确系统开发的技术路径与实施标准；

三是完善数据库相关内容，明确数据库表结构设计及初始化数据，保障数据存储安全、查询高效，支撑系统正常运行；

四是规范编制测试案例，全面覆盖核心功能测试与异常场景测试，明确测试步骤、预期结果及测试环境，为系统测试提供规范指引；

五是汇总测试执行情况、缺陷统计及测试结论，明确系统是否满足交付要求。

为大赛评审提供完整规范的参考依据，为项目后续开发、测试、部署及迭代优化提供全面指引，提升项目在大赛中的竞争力。

\subsection{项目概况}

\subsubsection{项目背景}

我国农业长期呈现“大国小农”格局，受从业人员不足、老龄化加剧、用地规模缩减等因素制约，借助高新技术革新传统生产模式、发展智慧农业成为转型必然。

智慧农业是现代信息技术与农业深度融合的高级形态，以信息和知识为核心，实现生产全流程感知、决策、控制、追溯的智能化，且我国智慧农业发展需推动系统国产化，掌握核心技术。叠加国家政策导向、市场需求升级，传统农业智能化管控不足、资源利用效率低等问题愈发突出。国家层面，战略与政策明确推动农业现代化转型，促进人工智能与农业发展相结合；行业层面，劳动力短缺与灌溉效率偏低的痛点，对智能化控制需求迫切；市场层面，智慧农业规模增速显著，全链路协同成为关注重点。

OpenHarmony的分布式核心特性，契合国产化、多端协同、低成本需求，是智慧农业设备协同的理想技术底座。

据此，我们开发基于OpenHarmony的智慧农业控制系统，适配服务外包大赛企业题目要求，破解传统农业痛点，助力农业提质增效。

\subsubsection{项目定位}

本项目定位为一款基于OpenHarmony生态的轻量化、高可靠、可演示的智慧农业控制系统，聚焦“监测、决策、执行、回传”一体化农业控制场景，形成以OpenHarmony应用端HAP为主控制入口、微信小程序为轻量运维入口、数字孪生平台为运行状态可视化入口的多端协同体系，实现“感知-分析-控制-追溯”全链路闭环管理。相较于行业主流方案，本项目依托OpenHarmony分布式特性实现多端与设备的统一开发，打通全控制链路并实现低成本移动部署，更贴合中小规模农业场景。

\subsubsection{项目目标}

\begin{itemize}[leftmargin=2em,itemsep=0.2em,topsep=0.2em]
    \item \textbf{功能目标：}实现环境数据（温湿度、光照等）实时采集与多端同步展示，智能灌溉、补光的自动与手动控制，演示/实物双模式切换，扩展设备接入及历史数据追溯等核心功能，满足任务要求；
    \item \textbf{性能目标：}确保数据传输延迟≤1秒、系统响应时间≤3秒，多端同步稳定，数据库读写高效，无明显卡顿，保障演示流畅性；
    \item \textbf{可靠性目标：}具备较强的容错能力，可应对设备离线、通信中断、数据异常等常见问题，系统运行稳定，无崩溃现象；
    \item \textbf{大赛目标：}完成系统全功能开发与测试，形成完整的提交材料，确保系统可正常演示，凸显技术创新性与实用性，提升大赛竞争力。
\end{itemize}

\section{需求分析}

\subsection{用户需求}

本系统核心面向农业种植户、系统维护人员两类主要使用者，两类用户共享统一后端数据模型，当前面向单一农业场景暂不设置独立权限分级，但界面入口与操作侧重不同，具体需求如下：
\begin{itemize}[leftmargin=2em,itemsep=0.2em,topsep=0.2em]

    \item \textbf{农业种植户：}核心需求是通过HAP端或小程序端完成日常农业生产远程管控，重点关注操作便捷性与状态反馈的及时性，无需复杂操作即可实现精细化管理，典型场景及具体需求包括：日常通过小程序快速查看温室大棚温湿度、光照等环境数据；根据系统推荐的灌溉水量，确认并执行灌溉操作；在异常天气等特殊场景下，远程调整补光参数，保障作物生长；可随时查询历史环境数据、灌溉及补光记录，辅助生产决策；界面以看板展示和一键控制为主，直观呈现核心信息，降低操作门槛。

    \item \textbf{系统维护人员：}核心需求是保障系统正常运行、完成设备与参数配置及故障排查，重点关注系统的可配置性与运行可追溯性，具体需求及典型场景包括：新传感器部署后，完成网关连接参数配置，确保设备正常入网；根据作物不同生长阶段，调整灌溉因子权重等策略阈值，适配种植需求；定期审查控制日志、操作日志，排查通信链路异常等问题，及时处理设备故障；负责系统日常运维，查看设备链路状态，确保数据采集、指令传输顺畅；界面以配置管理和日志查询为主，便捷开展运维操作，提升运维效率。

\end{itemize}

\subsection{功能需求}

为实现项目目标，系统在实地生产、远程运维和教学演示等应用场景中需覆盖以下核心功能：
\begin{itemize}[leftmargin=2em,itemsep=0.2em,topsep=0.2em]
  \item \textbf{环境感知与状态回传：}支持温度、湿度、光照、二氧化碳及设备状态持续采集与回传；前后端消息按统一业务类型分发并采用换行分隔，确保多端解析一致与状态展示连续。
  \item \textbf{智能灌溉联动控制：}基于温度、湿度、光照、历史灌溉、季节五因子计算推荐水量，推荐值限制在100\textasciitilde2000ml；执行阶段按5秒间隔回传进度、已用时和剩余时长，结束后回传100\%完成状态。
  \item \textbf{智能补光稳态控制：}按时段和季节修正计算目标亮度，支持自动与手动模式切换；自动模式采用最近5次滑动窗口与5\%最小变化阈值，亮度输出限制在0\textasciitilde100，降低抖动与频繁跳变。
  \item \textbf{多端协同控制：}HAP端、小程序端与数字孪生端共用统一的状态数据格式，\textbf{任一端下发请求均由后端分发处理并回传状态}，实现控制可达、状态可见、过程可追溯。
  \item \textbf{扩展设备接入：}系统在现有传感器与执行器基础上预留统一外设接入接口。以风扇控制为例，新增设备可通过在后端服务层增加对应控制逻辑，并在HAP端与控制页面补充启停组件，完成指令下发与状态回传全流程。
  \item \textbf{演示模式复现：}支持启停演示模式，在无实物硬件条件下持续生成并广播环境数据；仿真服务按0.5秒间隔刷新并广播数据、按1秒间隔写入数据库，保证演示连续性与可追溯性。
  \item \textbf{历史追溯：}自动记录实物小车端状态、控制命令、灌溉与补光状态，支持按时间范围和类型查询。
\end{itemize}

\subsection{非功能需求}

除功能需求外，系统还需在可靠性、实时性、可维护性和资源约束等方面满足以下要求：

可靠性方面，系统需实现端到端的链路保障机制。连接验证采用请求-应答模式，单次超时3秒、最大重试3次；后端到实物小车端维持3秒周期心跳，HAP端到后端维持5秒周期心跳；链路中断后支持自动重连与本地兜底策略，降低弱网与串口干扰导致的控制中断风险。

实时性方面，环境数据采集与回传需保证秒级刷新，仿真服务按0.5秒间隔广播实时数据；灌溉执行进度按5秒间隔推送，补光策略响应需在收到状态更新后的下一个计算周期内完成亮度调整；控制指令从下发到状态回传的端到端延迟应控制在可接受范围内，避免因延迟积累导致控制决策与实际状态脱节。

可维护性与可扩展性方面，系统采用模块化架构，后端服务按职责拆分为通信层、业务逻辑层和数据持久化层，新增设备仅需在对应层级扩展驱动与控制逻辑；关键流程均具备日志留痕能力，支持按时间范围和类型进行问题定位与运维审计。

资源约束方面，HAP端与数字孪生端历史序列上限均为60条，避免长期运行导致内存持续增长；传感器数据更新采用节流机制，在保证展示连续性的前提下平衡渲染性能与数据量。

\subsection{需求总结}

本文档需求分析围绕系统两类核心用户展开，全面梳理了用户需求、功能需求及非功能需求，形成完整的需求体系，为后续系统设计、开发及测试提供明确依据。

用户需求上，明确农业种植户便捷管控、系统维护人员高效运维的核心诉求，区分两类用户的操作侧重与场景需求；功能需求上，覆盖环境感知、智能控制、多端协同、双模式切换等核心模块，满足大赛任务与实际应用需求；非功能需求上，从性能、可靠性、可扩展性等维度提出明确标准，保障系统稳定运行与灵活适配。

整体需求紧扣项目定位，聚焦智慧农业管控痛点，兼顾大赛演示与实际落地，确保需求逻辑清晰、覆盖全面、可落地、可验证。

\subsection{项目价值}

在上述目标与需求基础上，本项目的核心价值体现在三方面。在业务层面，通过“采集—分析—控制—回传”的联动流程提升农业管理效率，降低人工巡检与经验决策成本。在技术层面，形成可复用的云边端协同架构与策略控制能力，为后续设备扩展和场景迁移提供基础。在推广层面，方案面向中小规模农业主体具备较强落地性，同时可作为教学示范的低成本智慧农业样板。

\section{系统设计}

\subsection{系统整体架构}

本系统采用四层架构设计，从下至上分别为硬件执行层、网关层、后端服务层、用户终端层，如图\ref{fig:system-architecture}所示。

\begin{figure}[H]
    \centering
    \includegraphics[width=0.95\linewidth]{pic/system-architecture.png}
    \caption{系统总体架构图}
    \label{fig:system-architecture}
\end{figure}

各层职责如下：
\begin{itemize}[leftmargin=2em,itemsep=0.2em,topsep=0.2em]
  \item \textbf{用户终端层：}OpenHarmony应用端HAP、微信小程序与数字孪生平台协同接入，分别承担实操控制、轻量运维和态势展示。
  \item \textbf{后端服务层：}Python后端承担协议解析、数据存储、策略计算、日志留痕与消息分发。
  \item \textbf{边缘网关层：}Hi3861承担UDP与UART桥接、命令下发、状态回传和链路保持。
  \item \textbf{硬件执行层：}STM32电机板与传感器板集成在轮式小车载体上，执行运动控制、超声测距与巡线检测等动作；温湿度与光照由Hi3861侧板载传感器采集，并支持按任务在棚区内移动部署。
\end{itemize}

\subsection{硬件系统架构设计}

本系统硬件架构对应总体架构中的“边缘网关层”与“硬件执行层”，整体采用“边缘网关层-硬件执行层”两层架构，各层设备协同工作，实现数据采集、动作执行与指令传输，兼顾部署成本与运行可靠性，适配实际移动部署需求，具体设计如下：
\begin{enumerate}[leftmargin=2em,itemsep=0.2em,topsep=0.2em]

    \item \textbf{边缘网关层：}核心设备为Hi3861网关，承担UDP与UART协议桥接、控制命令下发、设备状态回传、链路保持及本地数据缓存等核心职责，是硬件与软件系统交互的核心枢纽。其具体功能包括：

    \begin{itemize}[leftmargin=2em,itemsep=0.2em,topsep=0.2em]
        \item 接收后端服务层下发的JSON格式控制指令，通过UART协议桥接转换后下发至硬件执行层设备；
        \item 接收硬件执行层上报的各类数据（含传感器数据、设备执行状态），汇总封装为统一JSON格式，通过UDP协议回传至后端服务层；
        \item 读取自身侧板载I2C传感器（SHT20温湿度传感器、AP3216C光照传感器）的温湿光数据，完成数据初步采集与预处理；
        \item 维持与后端服务层、硬件执行层的稳定通信，确保链路通畅，应对通信中断等异常场景。
    \end{itemize}
    
    
    \item \textbf{硬件执行层：}核心设备集成在轮式小车载体上，包括STM32电机板、传感器板及各类执行器、传感器，主要负责数据采集（非网关侧板载传感器）与动作执行，具体组成及职责如下：

    \begin{itemize}[leftmargin=2em,itemsep=0.2em,topsep=0.2em]
        \item \textbf{STM32电机板：}作为硬件执行层核心控制单元，接收Hi3861网关下发的控制指令，执行运动控制、超声测距、巡线检测等动作；汇总传感器板采集的超声数据、巡线数据，将其封装为JSON格式，通过UART协议上报至Hi3861网关；联动灌溉、补光等执行设备，执行灌溉水量控制、补光亮度调节等操作，并将执行状态实时回传。

        \item \textbf{传感器板：}搭载超声传感器、巡线传感器，采集超声测距数据、巡线检测数据，通过UART二进制帧上报至STM32电机板，为轮式小车运动控制提供数据支撑；支持按任务在棚区内移动部署，提升数据采集的灵活性，适配中小规模农业种植场景的多样化需求。

        \item  \textbf{执行器与辅助传感器：}包括灌溉设备、补光设备，可扩展接入风扇等设备，接收STM32电机板的控制指令，执行相应操作；Hi3861网关侧板载SHT20温湿度传感器、AP3216C光照传感器，负责采集温湿度、光照强度数据，无需额外部署独立传感器，降低部署成本，数据直接由网关读取并上报后端。
    \end{itemize}
    
\end{enumerate}

总的来说，硬件架构各层交互为以下逻辑。硬件执行层的传感器板采集超声、巡线数据，通过UART二进制帧上报至STM32电机板；STM32电机板汇总数据并转换为JSON格式，经UART上报至Hi3861网关；Hi3861网关读取自身侧板载传感器数据，与STM32上报的数据整合封装，通过UDP协议回传至后端服务层；后端服务层下发的控制指令经UDP传输至Hi3861网关，网关通过UART协议桥接下发至STM32电机板，由电机板控制执行器完成相应动作，形成“采集-上报-指令-执行”的闭环链路。


\subsection{软件系统架构设计}

本系统软件架构对应总体架构中的“用户终端层”与“后端服务层”，整体采用“用户终端层-后端服务层”两层架构，与硬件架构的“边缘网关层-硬件执行层”协同联动，形成完整的四层总体架构体系。软件系统基于OpenHarmony分布式架构，按要求分为概要设计与详细设计两部分，确保各层职责清晰、交互顺畅。

\subsubsection{概要设计}

软件系统概要设计核心是明确软件两层架构的整体定位、模块划分及各模块核心职责，搭建软件系统整体框架，严格契合总体架构中用户终端层、后端服务层的职能要求，确保与边缘网关层、硬件执行层协同适配，满足需求分析中明确的功能与非功能要求，同时体现轻量化设计理念，以低成本、低复杂度实现核心功能。
\begin{enumerate}[leftmargin=2em,itemsep=0.2em,topsep=0.2em]
    \item \textbf{整体架构梳理：}明确软件系统“用户终端层-后端服务层”两层架构的核心定位，界定各层交互逻辑及与硬件架构的协同关系——用户终端层负责用户交互、实操控制、轻量运维与态势展示，发起控制请求并接收状态反馈；后端服务层采用Python后端，承担协议解析、数据存储、策略计算、日志留痕与消息分发核心职责，是软件系统的核心枢纽；后端服务层与边缘网关层（Hi3861网关）通过UDP协议实现数据交互，边缘网关层与硬件执行层通过UART协议实现数据交互，四层协同实现“采集-传输-处理-控制-反馈”全链路闭环。

    \item \textbf{核心模块划分：}按功能需求及两层架构职能，拆解为6大核心模块，各模块独立封装、协同工作，与总体架构职责精准对应，具体划分及核心职责如下：
    \begin{itemize}[leftmargin=2em,itemsep=0.2em,topsep=0.2em]

        \item \textbf{客户端交互模块：}归属用户终端层，覆盖OpenHarmony应用端HAP、微信小程序、数字孪生平台，分别承担实操控制、轻量运维和态势展示职能，与总体架构中用户终端层职责完全一致；负责环境数据可视化展示、控制指令下发、历史数据查询、模式切换等用户交互操作，核心是保障操作便捷性与多端同步一致性，支持多端协同接入，实现状态实时同步。
    
        \item \textbf{消息接收与分发模块：}归属后端服务层，作为后端统一消息入口，负责接收用户终端层通过UDP协议发送的JSON格式控制请求，以及边缘网关层通过UDP协议回传的JSON格式状态数据；根据消息类型（控制请求、状态数据）进行分类分发，将控制请求转发至对应业务模块，将状态数据转发至数据存储与追溯模块及客户端交互模块，确保消息流转高效、精准。
    
        \item \textbf{业务策略计算模块：}归属后端服务层，核心负责各类控制策略的计算与处理，对应总体架构中后端服务层“策略计算”职能；内置自动控制算法，基于边缘网关层上报的环境数据（温湿度、光照），结合预设的作物生长阈值，计算灌溉水量（100-2000ml）、补光亮度（0\%-100\%）等控制参数；接收用户手动控制请求，解析控制参数，生成标准化控制指令，转发至协议转换与指令下发模块。
    
        \item \textbf{协议转换与指令下发模块：}归属后端服务层，负责协议适配与控制指令下发，对应总体架构中后端服务层“协议解析”职能；将业务策略计算模块生成的控制指令封装为JSON格式，通过UDP协议下发至边缘网关层（Hi3861网关）；接收边缘网关层回传的JSON格式状态数据，完成数据格式校验与解析，确保数据准确性，为后续处理提供支撑；配合边缘网关层完成UDP与UART协议的协同衔接，保障指令与数据的双向顺畅传输。
    
        \item \textbf{虚实双模式管理模块：}归属后端服务层，负责实物模式（与硬件设备联动）与演示模式（无硬件条件下模拟全流程）的切换，适配大赛演示需求；演示模式下，模拟传感器数据生成、设备控制指令执行反馈及数据流转全流程，无需关联硬件设备即可完成演示；两种模式切换时，自动保存当前状态，切换后无缝衔接，确保系统兼顾演示性与实用性。
    
        \item \textbf{数据存储与追溯模块：}归属后端服务层，负责各类数据的写入、查询、管理与追溯，对应总体架构中后端服务层“数据存储、日志留痕”职能；对接数据库层，存储环境数据（温湿度、光照、超声、巡线数据）、设备状态、灌溉/补光记录、操作日志、控制指令等数据；支持按时间范围、设备ID等条件查询历史数据，为用户决策、系统运维提供数据支撑，确保链路可追溯。
    
    \end{itemize}

\end{enumerate}

\subsubsection{详细设计}

详细设计在概要设计基础上，明确各模块的内部实现逻辑、接口定义、数据交互格式，严格贴合总体架构中各层交互要求，为系统开发提供具体可落地的技术规范，重点包含模块详细实现、接口定义两部分，确保与硬件架构、总体架构协同一致。

\begin{enumerate}
    \item 各模块详细实现
    \begin{itemize}[leftmargin=2em,itemsep=0.2em,topsep=0.2em]
        \item \textbf{客户端交互模块：}分为HAP应用、微信小程序、数字孪生平台三部分，数据直观，操作简单。

        OpenHarmony应用端HAP基于OpenHarmony API开发，作为核心实操入口，承担实操控制职能，实现控制指令下发、实时数据展示、模式切换、设备状态监控等核心功能，适配OpenHarmony终端，操作流程简洁，满足种植户、维护人员的实操需求；微信小程序功能同HAP应用，同时无需复杂安装与操作，适配移动运维场景，降低运维复杂度，凸显“轻量便捷”的核心优势；数字孪生平台承担态势展示职能，实现设备状态、环境数据、数据流转链路的可视化建模，实时同步设备运行状态、控制指令执行情况，支持场景模拟展示，直观呈现系统运行态势。
        
        三者共用统一数据交互接口，通过UDP协议与后端服务层通信，确保数据同步延迟≤1秒，实现多端状态实时一致。

        \item \textbf{消息接收与分发模块：}采用UDP通信监听机制，持续监听用户终端层与边缘网关层的消息传输；内置消息解析规则，快速识别JSON格式消息中的消息类型、数据内容，区分控制请求（如灌溉控制、补光控制、模式切换）与状态数据（如环境数据、设备状态、执行结果）；建立消息分发路由，将控制请求精准转发至业务策略计算模块，将状态数据同步转发至数据存储与追溯模块，同时向所有已连接的客户端广播状态数据，实现多端状态同步；内置消息容错机制，应对消息丢失、格式错误等异常，确保消息流转稳定可靠。

        \item \textbf{业务策略计算模块：}内置两种控制模式的核心逻辑，适配不同使用场景：自动模式下，实时接收边缘网关层上报的环境数据（温湿度、光照），结合预设的作物生长阈值，通过内置算法计算最优控制参数，如灌溉水量（100~2000ml）、补光亮度（0~100），自动生成控制指令，无需人工干预；手动模式下，接收用户终端层下发的自定义控制参数，解析参数合法性（如灌溉水量、补光亮度是否在预设范围），验证通过后生成标准化控制指令；同时负责控制指令的状态管理，记录指令下发时间、执行状态，接收执行结果反馈，更新指令状态（待执行、执行中、已完成、失败），同步推送至客户端与日志模块。

        \item \textbf{协议转换与指令下发模块：}内置UDP与UART协议协同逻辑，配合边缘网关层完成协议桥接；将业务策略计算模块生成的控制指令封装为统一JSON格式，明确指令类型、设备ID、控制参数、操作人ID等核心信息，通过UDP协议下发至Hi3861网关，确保指令传输延迟≤1秒；接收Hi3861网关通过UDP回传的JSON格式状态数据，校验数据完整性与格式正确性，剔除错误数据，解析数据中的环境参数、设备状态、执行结果等核心信息，转发至消息接收与分发模块；支持协议适配扩展，可适配新增设备的通信协议，提升系统可扩展性。

        \item \textbf{虚实双模式管理模块：}内置模式切换开关，默认启动实物模式（与硬件执行层、边缘网关层联动），用户可通过用户终端层手动切换至演示模式；演示模式下，模拟传感器数据生成（按合理范围随机生成温湿度、光照、超声、巡线数据），模拟控制指令下发、执行反馈及数据流转全流程，无需关联硬件设备，即可完成系统全功能演示；两种模式切换时，自动保存当前系统状态（如当前环境数据、设备状态、控制参数），切换后无缝衔接，确保演示与实际运行场景的一致性；内置模式标识，同步至各模块与数据库，确保数据存储与展示的准确性。

        \item \textbf{数据存储与追溯模块：}对接数据库层，提供数据写入、查询、删除、更新等标准化接口；按时间戳对数据进行分区存储，将环境数据、控制记录、操作日志分类存储，提升查询效率；支持按时间范围、设备ID、数据类型（传感器数据、灌溉记录、补光记录）等条件查询历史数据，查询响应时间≤10秒；定期对数据进行备份，防止数据丢失，同时记录操作日志、控制日志，支撑系统运行可追溯，便于维护人员排查链路异常、设备故障；同步向客户端交互模块推送实时数据，支撑多端数据可视化展示。

    \end{itemize}

    \item \textbf{接口定义}
    
    本系统接口按“客户端层↔后台服务层”、“后台服务层↔Hi3861网关层”、“Hi3861网关层↔STM32电机板”、“STM32电机板↔STM32传感器板”四级划分，同时包含数据库服务接口，明确各层级通信方式、协议格式、消息类型及参数规范，适配系统轻量化设计与多端协同需求
    
\end{enumerate}

\subsection{通信与协议设计}

系统通信链路采用“上层JSON、底层二进制与本地I2C采集并行”的分层设计，兼顾可读性与实时性。全系统共涉及五条主要通信链路：HAP/小程序与后端之间、后端与Hi3861网关之间采用UDP传输JSON文本，保障多端交互与策略下发的一致性；Hi3861与STM32电机板之间采用UART串口传输JSON短键名；电机板与传感器板之间采用UART自定义二进制帧以保证超声与巡线采样实时性；温湿度与光照等环境数据由Hi3861通过I2C直连SHT20与AP3216C完成采集。各链路的具体参数配置汇总如表\ref{tab:communication-link-protocol}所示。本节描述的通信链路为后续灌溉、补光、移动控制等功能的指令下发与状态回传提供了统一的传输基础。

\begin{table}[htbp]
  \centering
  \caption{通信链路与协议设计表}
  \label{tab:communication-link-protocol}
  \setlength{\tabcolsep}{4pt}
  {\renewcommand{\arraystretch}{1.28}
  \begin{tabular}{L{0.36\linewidth}L{0.15\linewidth}L{0.16\linewidth}L{0.25\linewidth}}
    \hline
    链路 & 传输协议 & 数据格式 & 关键参数 \\
    \hline
    HAP/小程序 $\leftrightarrow$ 后端 & UDP & JSON文本 & 端口 8888 \\
    后端 $\leftrightarrow$ Hi3861 & UDP & JSON & 端口 7788 \\
    Hi3861 $\leftrightarrow$ STM32电机板 & UART & JSON & 波特率 115200 \\
    STM32电机板 $\leftrightarrow$ STM32传感器板 & UART & Binary & 波特率 115200 \\
    Hi3861 $\leftrightarrow$ 温湿度/光传感器 & I2C & 寄存器读写 & 温湿度/光照本地采集 \\
    \hline
  \end{tabular}
  }
\end{table}

\vspace{8pt}

由表\ref{tab:communication-link-protocol}可见，系统在业务交互链路采用JSON文本，在设备控制链路采用UART的JSON短键名与二进制帧，并在环境感知链路通过I2C本地采集温湿度与光照数据；Hi3861统一完成协议桥接与数据汇聚。

在上层JSON链路中，Hi3861与STM32电机板之间的UART通信采用JSON短键名格式以减少传输字节数。典型下行指令示例如：运动控制使用\texttt{\{"t":"move","d":"f","s":50\}}（t为类型、d为方向、s为速度百分比）；灌溉控制使用\texttt{\{"t":"water","v":500\}}（v为水量ml）；补光控制使用\texttt{\{"t":"light","v":60\}}（v为亮度百分比）。上行状态回传采用\texttt{\{"t":"status","spd":30,"bat":85,"dst":15.2\}}（速度、电量、距离）等短键名结构。相比上层完整键名，短键名将单条消息长度控制在30\textasciitilde50字节以内，在115200波特率下单帧传输时间约为2.6\textasciitilde4.3ms，满足运动控制的实时性要求。

在底层设备链路中，STM32电机板与传感器板之间的UART通信采用自定义二进制帧格式。帧结构由帧头（2字节固定标识0xAA55）、帧长（1字节）、命令字（1字节）、数据域（可变长度）和校验和（1字节异或校验）五部分组成。二进制帧的编码效率高于JSON文本——以超声测距数据为例，二进制帧仅需8字节即可传输左右传感器距离值（各2字节），而等效JSON文本至少需要25字节以上。电机板在收到二进制帧后通过帧头和校验和进行完整性验证，校验失败时丢弃当前帧并等待下一帧，避免解析错误导致的执行异常。

在可靠性方面，系统已实现心跳保活、接收超时、异常重置和状态回传机制，降低弱网与串口干扰造成的中断风险。

在通信协议选型的学术论证方面，智慧农业IoT领域主流的应用层协议包括MQTT、CoAP和HTTP等，各自的适用场景、性能指标及本项目的选型依据对比如表\ref{tab:protocol-comparison}所示。

\begin{table}[H]
  \centering
  \caption{应用层协议定量对比与选型依据}
  \label{tab:protocol-comparison}
  \setlength{\tabcolsep}{3pt}
  {\renewcommand{\arraystretch}{1.20}
  \begin{tabular}{L{0.12\linewidth}L{0.20\linewidth}L{0.19\linewidth}L{0.19\linewidth}L{0.20\linewidth}}
    \hline
    维度 & MQTT & CoAP & HTTP & 本项目（UDP+JSON） \\
    \hline
    传输层 & TCP & UDP & TCP & UDP \\
    连接开销 & 需建立Broker连接 & 无连接 & 需建立TCP连接 & 无连接 \\
    典型延迟 & 50\textasciitilde200ms & 10\textasciitilde50ms & 100\textasciitilde500ms & 1\textasciitilde10ms（局域网） \\
    最小帧头 & 2字节 & 4字节 & 数百字节 & 无固定帧头 \\
    可靠性机制 & QoS 0/1/2 & CON/NON消息 & TCP重传 & 应用层心跳+重试 \\
    适用场景 & 多设备大规模订阅 & 受限环境低功耗 & Web服务交互 & 局域网点对点控制 \\
    部署复杂度 & 需部署Broker & 需CoAP服务器 & 需HTTP服务器 & 无中间件依赖 \\
    \hline
  \end{tabular}
  }
\end{table}

由表\ref{tab:protocol-comparison}可见，本系统选择UDP传输JSON文本的方案基于以下学术与工程双重考量。在延迟层面，系统运行于局域网环境（后端与Hi3861处于同一网段），UDP无连接建立的握手开销，单次指令传输延迟典型值为1\textasciitilde10ms，相比MQTT的TCP三次握手和Broker转发可降低一个数量级，满足灌溉与补光控制的实时性要求。在可靠性层面，MQTT的QoS机制虽然提供从“至多一次”到“恰好一次”的交付保证，但其代价是增加Broker中间件的部署和维护开销；对于当前系统客户端数量有限（HAP端+小程序端+数字孪生端共3类）的局域网场景，网络丢包率通常低于0.1\%，通过应用层的心跳保活（5秒周期）和3次超时重试即可达到工程可接受的可靠性水平，无需引入MQTT Broker的额外中间层\cite{smart_agri_iot2025}。在部署复杂度层面，UDP方案无需部署额外的消息中间件，后端Python进程即可完成所有收发逻辑，降低了中小规模农业主体的运维门槛。需要指出的是，当系统未来扩展到大规模多设备场景（数十个以上网关并发接入、广域网跨地域通信）时，MQTT的发布/订阅模式和QoS保障将成为更优选择，当前架构已通过统一消息类型设计为后续协议迁移预留了接口。

在心跳保活与连接验证的状态机设计方面，HAP端的连接管理遵循“断开→验证中→已连接→保活中→断开”的状态转换流程。初始处于“断开”状态时，用户点击连接后发送\texttt{connect}请求并进入“验证中”状态，启动3秒超时计时器；若在超时前收到\texttt{connect\_ack}应答则转入“已连接”状态，否则在最多3次重试后返回“断开”状态并提示用户。进入“已连接”状态后，系统自动启动5秒周期心跳，每次发送\texttt{ping}消息后等待\texttt{pong}应答；若连续3次心跳未收到应答则判定链路中断，自动切回“断开”状态。后端到Hi3861网关侧采用类似的3秒周期心跳机制，但超时判定策略不同：若单次心跳超时则立即尝试重发初始化命令\texttt{carStatus:on}，连续超时3次后标记车端离线并停止向客户端转发该车端数据。

在 I2C 传感器采集方面，Hi3861 通过硬件 I2C 总线直连 SHT20 温湿度传感器和 AP3216C 光照传感器。SHT20 支持标准 I2C 通信（设备地址 0x40），温度和湿度分别通过写入不同的命令字触发测量，等待测量完成后读取 2 字节数据加 1 字节 CRC 校验。以温度为例，读取命令为 0xE3，返回数据通过公式 $T = -46.85 + 175.72 \cdot \frac{S_T}{2^{16}}$ 转换为实际温度值（$^\circ$C），其中 $S_T$ 为 16 位原始 ADC 值；湿度读取命令为 0xE5，转换公式为 $RH = -6 + 125 \cdot \frac{S_{RH}}{2^{16}}$（\%RH）。AP3216C（设备地址 0x1E）内置环境光传感器和接近传感器，光照数据通过读取 0x0C\textasciitilde0x0E 三个寄存器获得，分辨率可达 65536 级。Hi3861 在主循环中以固定周期触发传感器采集，将读取到的温湿光数据与设备状态封装为统一 JSON 格式上报后端，实现了“本地采集—本地组帧—远程上报”的轻量化感知链路\cite{smart_agri_sensor2025}。这些由底层硬件实时生成的封装报文构成了数据上行流转的起点，配合由用户侧自上而下发起的控制请求，构成了跨越终端、服务、网关与硬件层的数据交互路径，其具体的指令往返与状态回传时序如图\ref{fig:cloud-edge-end-path}所示。

\begin{figure}[H]
    \centering
    \includegraphics[width=1\linewidth]{pic/cloud-edge-end-path.png}
    \caption{云边端数据与指令往返路径}
    \label{fig:cloud-edge-end-path}
\end{figure}

在上述链路框架下，通信层面落实了以下七项策略，旨在保障异构终端间交互的确定性与数据流转的稳定性：
\begin{itemize}[leftmargin=2em,itemsep=0.2em,topsep=0.2em]
  \item \textbf{消息分发统一：}后端对客户端上报消息按业务类型集中分发，覆盖连接、控制、查询、演示模式、灌溉与补光等类别，其中连接验证消息优先处理。
  \item \textbf{网关桥接与链路维持：}Hi3861在上行侧承接后端UDP收发，在下行侧完成UART桥接与命令转发，并通过I2C直连传感器采集温湿度与光照数据，配合状态回传与重连流程保障通信连续性。
  \item \textbf{连接验证规范：}HAP端在发起连接时先发送验证请求，在规定时间内收到应答后才确认连接建立并启动心跳；超时则在上限范围内自动重试。
  \item \textbf{弱网容错与异常恢复：}链路侧启用心跳保活、连接重试、接收超时重置与本地兜底策略，降低弱网干扰下的数据丢包与时延波动风险。
  \item \textbf{双模式数据格式一致：}实物小车回传的环境数据与演示模式生成的仿真数据采用相同消息结构，前端界面无需区分来源即可统一渲染。
  \item \textbf{广播路径清晰：}实物模式下，数据仅向已连接且未开启演示的客户端推送；演示模式下，仿真数据通过独立广播通道同步到监听端。
  \item \textbf{过程数据可审计：}车端状态数据以秒为单位持续写入数据库，控制命令同步记录来源、类型与演示标记，为复盘与故障定位提供依据。
\end{itemize}

\subsection{数据库设计}

\subsubsection{总体设计}

本数据库覆盖环境监测、智能灌溉、智能补光、设备管理、指令控制、数据追溯、演示仿真全业务场景，满足系统可追溯审计、部署轻量化、虚实双模式分离等需求。数据库基本信息如表\ref{tab:数据库基本信息}所示。

\begin{table}[htbp]
  \centering
  \caption{数据库基本信息}
  \label{tab:数据库基本信息}
  \begin{tabular}{|l|l|}
    \hline
    项目     & 配置                  \\
    \hline
    数据库名 & smart\_car            \\
    字符集   & utf8mb4               \\
    排序规则 & utf8mb4\_unicode\_ci  \\
    存储引擎 & InnoDB                \\
    \hline
  \end{tabular}
\end{table}

\subsubsection{数据表详细结构}

\paragraph{3.5.2.1实际传感器数据表}

实际传感器数据表（如表\ref{tab:实际传感器数据表}所示）用于存储实物模式下，Hi3861 采集的真实环境与设备状态（1 秒 / 条）。

\begin{table}[htbp]
  \centering
  \caption{实际传感器数据表 car\_sensor\_data}
  \label{tab:实际传感器数据表}
  \tiny
  \begin{tabular}{|l|l|l|l|}
    \hline
    字段名          & 类型         & 约束                        & 说明              \\
    \hline
    id              & BIGINT       & PK, AUTO\_INC               & 自增ID            \\
    device\_id       & VARCHAR(50)  & NOT NULL, INDEX             & 设备ID            \\
    timestamp       & DATETIME(3)  & NOT NULL, INDEX             & 采集时间          \\
    temperature     & DECIMAL(5,2) &                             & 实际温度(℃)       \\
    humidity        & DECIMAL(5,2) &                             & 实际湿度(\%)       \\
    lux             & INT          &                             & 实际光照强度(lux) \\
    proximity       & INT          &                             & 接近距离          \\
    ir\_value        & INT          &                             & 人体检测值        \\
    co2             & INT          &                             & 模拟CO2浓度(ppm)  \\
    car\_status      & VARCHAR(20)  &                             & 小车状态          \\
    car\_mode        & VARCHAR(20)  &                             & 运行模式          \\
    left\_speed      & INT          &                             & 左电机速度(mm/s)  \\
    right\_speed     & INT          &                             & 右电机速度(mm/s)  \\
    battery\_voltage & INT          &                             & 电池电压(mV)      \\
    distance        & INT          &                             & 障碍物距离(mm)    \\
    fan             & TINYINT      &                             & 风扇状态(0/1)     \\
    led             & TINYINT      &                             & LED状态(0/1)      \\
    buzzer          & TINYINT      &                             & 蜂鸣器状态(0/1)   \\
    created\_at      & DATETIME     & DEFAULT CURRENT\_TIMESTAMP  & 创建时间          \\
    \hline
  \end{tabular}
\end{table}

\paragraph{3.5.2.2模拟传感器数据表}

模拟传感器数据表（如表\ref{tab:模拟传感器数据表}所示）用于存储演示模式下后端生成的仿真环境数据，在无硬件情况下也可完整演示。

\begin{table}[htbp]
  \centering
  \caption{模拟传感器数据表 simulated\_sensor\_data}
  \label{tab:模拟传感器数据表}
  \begin{tabular}{|l|l|l|l|}
    \hline
    字段名      & 类型         & 约束                        & 说明              \\
    \hline
    id          & BIGINT       & PK, AUTO\_INC               & 自增ID            \\
    device\_id   & VARCHAR(50)  & NOT NULL, INDEX             & 设备ID            \\
    timestamp   & DATETIME(3)  & NOT NULL, INDEX             & 模拟时间          \\
    temperature & DECIMAL(5,2) &                             & 模拟温度(℃)       \\
    humidity    & DECIMAL(5,2) &                             & 模拟湿度(\%)       \\
    lux         & INT          &                             & 模拟光照强度(lux) \\
    co2         & INT          &                             & 模拟CO2浓度(ppm)  \\
    time\_period & INT         &                             & 时间段(0-6)       \\
    created\_at  & DATETIME    & DEFAULT CURRENT\_TIMESTAMP  & 创建时间          \\
    \hline
  \end{tabular}
\end{table}

\paragraph{3.5.2.3控制指令表}

控制指令表（如表\ref{tab:控制指令表}所示）用于存储所有来自 HAP、微信小程序、数字孪生的控制指令，在必要时用于审计，可追溯、可复盘。

\begin{table}[htbp]
  \centering
  \caption{控制指令表 control\_commands}
  \label{tab:控制指令表}
  \begin{tabular}{|l|l|l|l|}
    \hline
    字段名       & 类型        & 约束                        & 说明               \\
    \hline
    id           & BIGINT      & PK, AUTO\_INC               & 自增ID             \\
    device\_id    & VARCHAR(50) & NOT NULL, INDEX             & 设备ID             \\
    timestamp    & DATETIME(3) & NOT NULL, INDEX             & 指令发送时间       \\
    command\_type & VARCHAR(50) & NOT NULL                    & 指令类型           \\
    command\_data & JSON        &                             & 指令详细内容       \\
    source       & VARCHAR(20) & DEFAULT 'miniapp'           & 指令来源           \\
    is\_simulated & TINYINT     & DEFAULT 0                   & 是否演示模式下发送 \\
    created\_at   & DATETIME    & DEFAULT CURRENT\_TIMESTAMP  & 创建时间           \\
    \hline
  \end{tabular}
\end{table}

\paragraph{3.5.2.4智能灌溉记录表}

智能灌溉记录表分为灌溉配置表和灌溉历史表（分别如表\ref{tab:灌溉配置表}所示和如表\ref{tab:灌溉历史表}所示）分别用于记录五因子算法灌溉执行全过程，包含推荐值、状态、进度、因子明细。

\begin{table}[htbp]
  \centering
  \caption{灌溉配置表irrigation\_config}
  \label{tab:灌溉配置表}
  \begin{tabular}{|l|l|l|l|}
    \hline
    字段名                  & 类型        & 约束                                   & 说明             \\
    \hline
    id                      & INT         & PK, AUTO\_INC                          & 自增ID           \\
    device\_id               & VARCHAR(50) & UNIQUE, NOT NULL, INDEX                & 设备ID           \\
    base\_water\_amount       & FLOAT       & DEFAULT 500                            & 基础灌溉水量(ml) \\
    humidity\_threshold      & FLOAT       & DEFAULT 40                             & 空气湿度阈值(\%)  \\
    auto\_irrigation\_enabled & TINYINT     & DEFAULT 0                              & 是否启用自动灌溉 \\
    scheduled\_time          & TIME        &                                        & 定时灌溉时间     \\
    created\_at              & DATETIME    & DEFAULT CURRENT\_TIMESTAMP             & 创建时间         \\
    updated\_at              & DATETIME    & DEFAULT CURRENT\_TIMESTAMP ON UPDATE    & 更新时间         \\
    \hline
  \end{tabular}
\end{table}

\begin{table}[htbp]
  \centering
  \caption{灌溉历史表 irrigation\_history}
  \label{tab:灌溉历史表}
  \begin{tabular}{|l|l|l|l|}
    \hline
    字段名            & 类型        & 约束                        & 说明                   \\
    \hline
    id                & INT         & PK, AUTO\_INC               & 自增ID                 \\
    device\_id         & VARCHAR(50) & NOT NULL, INDEX             & 设备ID                 \\
    start\_time        & DATETIME    & NOT NULL                    & 灌溉开始时间           \\
    end\_time          & DATETIME    &                             & 灌溉结束时间           \\
    water\_amount      & FLOAT       & NOT NULL                    & 灌溉水量(ml)           \\
    duration          & INT         &                             & 灌溉时长(秒)           \\
    status            & ENUM        & DEFAULT 'pending'           & 灌溉状态               \\
    trigger\_type      & ENUM        & DEFAULT 'manual'            & 触发类型               \\
    sensor\_data       & JSON        &                             & 灌溉时传感器数据快照   \\
    algorithm\_factors & JSON        &                             & 算法各因子数值         \\
    created\_at        & DATETIME    & DEFAULT CURRENT\_TIMESTAMP  & 创建时间               \\
    \hline
  \end{tabular}
\end{table}


\paragraph{3.5.2.5智能补光状态表}

智能补光状态表（如表\ref{tab:智能补光状态表}所示）用于记录补光策略执行、模式切换、亮度输出，支持自动 / 手动双模式追溯。

\begin{table}[htbp]
  \centering
  \caption{智能光照状态表 smart\_light\_status}
  \label{tab:智能补光状态表}
  \begin{tabular}{|l|l|l|l|}
    \hline
    字段名            & 类型        & 约束                        & 说明                   \\
    \hline
    id                & BIGINT      & PK, AUTO\_INC               & 自增ID                 \\
    device\_id         & VARCHAR(50) & NOT NULL, INDEX             & 设备ID                 \\
    timestamp         & DATETIME(3) & NOT NULL, INDEX             & 时间戳                 \\
    mode              & TINYINT     &                             & 模式(1=auto, 0=manual) \\
    brightness        & INT         &                             & 当前亮度(0-100)        \\
    target\_brightness & INT         &                             & 目标亮度               \\
    time\_period       & INT         &                             & 时间段                 \\
    light\_level       & INT         &                             & 光照等级(0-5)          \\
    created\_at        & DATETIME    & DEFAULT CURRENT\_TIMESTAMP  & 创建时间               \\
    \hline
  \end{tabular}
\end{table}

\subsubsection{数据库使用规范}

数据保留：传感器数据默认保留90 天，审计 / 记录数据永久保留。

写入频率：实物数据 1 秒 / 条；演示数据 0.5 秒刷新、1 秒入库。

查询优化：所有时间范围查询必须走device\_id + ts联合索引。

扩展兼容：新增设备（风扇、卷帘机）只需在sensor\_data\_real扩展字段，无需改表结构。

\section{测试案例}

\subsection{测试概述}

\subsection{核心功能测试案例}

\subsection{异常场景测试案例}

\subsection{系统性能测试案例}

\section{测试报告}

\subsection{测试概况}
\subsection{测试内容}

\subsubsection{核心功能测试内容}

\subsubsection{异常场景测试内容}

\subsubsection{系统性能测试内容}

\subsection{测试结果汇总}
\subsection{测试结果分析}
\subsubsection{功能测试结果分析}
\subsubsection{异常测试结果分析}
\subsubsection{性能测试结果分析}

\end{document}
